% !TeX root = ../main.tex

\begin{framed}\large\bfseries
  \#2 Free Electron Model of Metals
\end{framed}

\paragraph{Drude model}\leavevmode

\subparagraph{Basic assumptions}\leavevmode

\begin{enumext*}[nosep, columns = 5, label = (\alph*)]
  \item(5) Free electron approximation
  \item(5) Electrical resistance results from scattering
  (all information is lost upon scattering).
  \item(3) Relaxation-time approximation.
  \item(2) Thermal equilibrium.
\end{enumext*}

\subparagraph{With(out) electric field}

Free electron moving randomly, no net current;
Or the drift velocity $v_d \ll v_\text{rms}$.

\subparagraph{EOM}

\begin{enumext}[nosep]
  \item Let $\tau$ be the scattering time and $1/\tau$ be the scattering rate.
  \item In $\d t$, a fraction $\frac{\d t}\tau$ of the electrons are scattered
  with zero average velocity.
  \item The rest part $1 - \d t/\tau$ are accelerated by external force $-eE$.
  \item In $t \sim t + \d t$, the EOM (omit high-order terms)
  \[
    m\odv{\bm v}/t = -m\bm v/\tau - e\bm E
  \]
  At steady state ($\odv{\bm v}t = 0$). The drift velocity
  \[
    \bm v_d \equiv \braket<\bm v(t)> = -e\tau\bm E/m \equiv -\mu\bm E,
  \]
  where mobility $\mu = \frac{\bm v}{\bm E} = \frac{e\tau}{m}$.
  For a given field, a higher $\mu$ means the electrons will move faster.
  The conductance
  \[
    \sigma = \bm j/\bm E
  = -en\bm v/\bm E = ne\mu = -e^2n\tau/m^*.
  \]
  Here, $\tau = l/v_\text{rms}$ is usually independent of $E$.
\end{enumext}

\subparagraph{Classical Hall effect}

\begin{enumext}
  \item Classically, electrons experience
  $\bm F = -e(\bm E + \bm v \times \bm B)$
  \item Cyclotron motion: $\bm E = 0$
  \[
    \ddot x + \ab(eB_z/m^*)^2x = 0, \enspace
    \omega_c = eB/m^*, \enspace
    r_c = v/\omega_c \propto 1/B.
  \]
  Classically, any orbit size is allowed.
\end{enumext}

\subparagraph{Drude model with a magnetic field}\leavevmode

Consider Drude model with applied electric and magnetic fields
\[
  m^*\dot{\bm v} = -e(\bm E + \bm v \times \bm B) - m^*\bm v/\tau
\]
Taking the steady state, with Ohm's rule $\bm j = -ne\bm v = \hat\sigma \bm E$
\begin{gather*}
  \hat\sigma =
  \begin{pmatrix}
    \frac{\sigma_0}{1 + \omega_c^2\tau^2} &
    \frac{\sigma_0\omega_c\tau}{1 + \omega_c^2\tau^2} & 0\\
   -\frac{\sigma_0\omega_c\tau}{1 + \omega_c^2\tau^2} &
    \frac{\sigma_0}{1 + \omega_c^2\tau^2} & 0\\
    0 & 0 & \sigma_0
  \end{pmatrix},\
  \hat\rho = \hat\sigma^{-1},\
  \sigma_0 = \frac{ne^2\tau}{m^*}.
\end{gather*}
The Hall coefficient $R_H = \rho_{xy}/B = -1/ne$.

\subparagraph{Drawbacks}

\begin{enumext*}[label = (\alph*), nosep, columns = 7]
  \item(3) Existence of insulators.
  \item(4) Temperature dependence of $\sigma$.
  \item(7) Heat capacity $C_v^\text D = \frac32nk_B$ does not match
  $C_v \approx 10^{-4} nk_BT$.
\end{enumext*}

\paragraph{Summerfeld model for electrons in metals}

\begin{enumext}
  \item Electrons treated quantum mechanically:
  Quantize the energy of electron gas.
  \item The electronic velocity distribution is given by the quantum
  Fermi-Dirac distribution.
  \item Pauli exclusion: At most one electron can occupy any single
  electron level.
\end{enumext}

\subparagraph{Electron wave}\leavevmode

Schr\"odinger equation for free electron
$\nabla^2\psi(\bm r) = -\frac{2mE}{\hbar^2}\psi(\bm r)$.
The solution is a plane wave with the energy
\[
  \psi(\bm r) = \frac1{\sqrt V} \upe^{\iu(k_xx + k_yy + k_zz)},\enspace
  E = \frac{\hbar^2(k_x^2 + k_y^2 + k_z^2)}{2m}.
\]

\subparagraph{Born-von Karman periodic boundary conditions}\leavevmode

The electrons are confined in a 3D box ($V = L_xL_yL_z$) with
\begin{enumext}[columns = 2]
  \item $U(\bm r) = 0$ inside the box,
  \item $U(\bm r) = \infty$ outside the box.
\end{enumext}
The PBC can be applied
\[
  \psi(x + L_x, y, z) = \psi(x, y + L_y, z) = \psi(x, y, z + L_z)
= \psi(x, y, z),
\]
determines the allowed $k$
\[
  \upe^{\iu k_iL_i} = 1,\quad k_i = 2\pi n_i/L_i, \quad n_i \in \mathbb Z.
\]
The boundary conditions lead to quantization.

\subparagraph{States in $k$-space}\leavevmode

\begin{enumext}
  \item The allowed quantum states can be visualized as a 3D grid of points in
  $k$-space with the distance of $\Delta k_i = 2\pi/L_i$.
  \item In $k$-space, one grid point occupies
  $\prod_{x,y,z}(\frac{2\pi}{L_i}) = (2\pi)^3/V$.
  \item There are $2 \times V/(2\pi)^3$ (Each grid point can be occupied
  by $\ket|\uparrow>$, $\ket|\downarrow>$) quantum states per unit
  volume in k-space.
\end{enumext}

\subparagraph{Fermi energy}\leavevmode

\begin{enumext}
  \item The free electron gas is isotropic, so the surface of constant
  $E = \hbar^2(k_x^2 + k_y^2 + k_z^2)/2m$ in k-space is a sphere.
  \item Fermi energy: energy of the highest occupied state at zero temperature.
  \begin{enumext}
    \item Number of
    electrons: $N = 2 \times \frac{4\pi k_F^3/3}{(2\pi/L)^3} = k_F^3V/3\pi^2$.
    \item Fermi wave vector: $k_F = (3\pi^2n)^{1/3}$, $n = N/V$ denotes the
    carrier density.
    \item Fermi
    energy: $E_f = \hbar^2k_F^2/2m = \frac{\hbar^2}{2m}(3\pi^2n)^{2/3}$.
    \item Fermi velocity and temperature relate to the Fermi
    energy $E_F = \frac12mv_F^2 = k_BT_F$.
  \end{enumext}
\end{enumext}

\subparagraph{Bosons and Fermions at $T = 0$}

\begin{enumext}
  \item Bosons: The occupation number can take on any value:
  photons, mesons, any composite particles with integer spin.
  \item Fermions: The occupation number can be zero or one due to the Pauli
  exclusion principle: lectrons, protons, neutrons, any composite
  particles with half-integer spin.
\end{enumext}

\subparagraph{Bosons and Fermions at $T \neq 0$}\leavevmode

At finite temperature, Bosons/Fermions obeys Bose-Einstein/Fermi-Dirac
statistics $\bar n_i = g_i\{\exp[(\epsilon_i-\mu)/k_BT] \mp 1\}^{-1}$.
\begin{enumext}[columns = 2]
  \item $\epsilon_i$: Energy of the $i$-th state.
  \item $n_i$: Occupation number.
  \item $g_i$: Degeneracy.
  \item $\mu$: Chemical potential.
\end{enumext}
Bosons and Fermions are identical and indistinguishable particles

\subparagraph{The electron gas at zero temperature}

\begin{enumext}
  \item  All quantum sates inside the Fermi sphere are filled while outside
  are empty.
  \item The surface of the Fermi sphere is called Fermi surface.
\end{enumext}

\subparagraph{Fermi-Dirac distribution}\leavevmode

Probability distribution function $f(E)$: the probability that a state
with energy $E$ is occupied.
\[
  f(E) = [1 + \upe^{(E-E_f)/k_BT}]^{-1} \xlongrightarrow{T=0}
  \begin{cases}
    0, & E > E_f,\\
    1, & E \leq E_f.
  \end{cases}
\]
As $T$ increasing, the function will broaden $\sim 2k_BT$ at $E = E_f$.
\begin{enumext}
  \item At Fermi level, the probability being occupied is always $1/2$.
  \item $E_f(T = 0) = E_F$, $E_f(T > 0) < E_F$
\end{enumext}

\subparagraph{Density of states}

The number states can be occupied per energy.
\[
  g(E) = \odv nE = \frac1V\odv NE =
  \begin{cases*}
    \frac{m^*}{\pi^2\hbar^3}\sqrt{2m^*E}, & 3D,\\
    \frac1{\pi\hbar} \sqrt{2m^*/E},       & 1D,\\
  \end{cases*}\
  \begin{array}{ll}
    m^*/\pi\hbar^2,                       & 2D,\\
    2\delta(E),                           & 0D.
  \end{array}
\]

\subparagraph{The electron gas at $T > 0$}

\begin{enumext}
  \item The total number of electron at $T > 0$
  \[
    n = N/V = \int_0^\infty g(E) [\upe^{(E-E_f)/k_BT} + 1]^{-1} \d E.
  \]
  \item  The energy
  density $u = U/V = \int_0^\infty Eg(E) \frac1{\upe^{(E-E_f)/k_BT} + 1} \d E$.
\end{enumext}

\subparagraph{Current density}

\[
  \bm j =-ne\bm v = -2e \int \frac{\d^3\bm k}{(2\pi)^3} f(\bm k) v(\bm k).
\]
At zero field, for every occupied state $\bm k$, there is a state $-\bm k$
occupied with the same probability. Therefore, $\bm j = 0$.

In Drude model, assume that scattering adds damping
\[
  \odv{\hbar\bm k(t)}t = -e\bm E - \frac{\hbar[\bm k(t) - \bm k]}{\tau}
  \xlongrightarrow[\text{steady state}]{t=\infty} \bm k(t)
= \bm k - \frac{e\tau}{\hbar}\bm E,
\]
the wave vector of every electron is shifted. The current density
\[
  \bm j = -2e\int \frac{\d^3\bm k}{(2\pi)^3}
          f\ab(\bm k + \frac{e\tau}\hbar\bm E) v(\bm k)
        \xlongequal{\int\d^3\bm k f(\bm k)\bm k = 0}
        \frac{e^2\tau}{m}n\bm E
\]
with $\bm j = \sigma\bm E$, $\sigma = ne^2\tau/m$ the same as that in
Drude model.

The excess current carriers are only those very near the Fermi surface.
So the current carriers have velocity $v_F$.

\subparagraph{Heat capacity of an electron gas}

At low temperature, $\mu \sim E_f$ and $\pdv{f(E,T)}T$ is peaked close to $E_f$,
$k_BT \ll E_f$. Then
\[
  C_e = \pdv{\Delta U}/T \approx \pi^2 nk_B^2T/2T_f \propto T.
\]
The heat capacity of electrons increases linearly with temperature.

\subparagraph{Limitation of Free electron gas model}

\begin{enumext}
  \item Cannot explain semiconductors/insulators.
  \item $C_V \propto T$ only at very low $T$.
  \item Fail to explain $C_V \sim T^3$ at intermediate temperature region.
  \item Wiedemann-Franz law temperature dependent.
  \item The sign of Hall coefficient $R_H = -1/ne$.
\end{enumext}

\begin{framed}\large\bfseries
  \#3 Lattice Vibrations and Phonons
\end{framed}

\subparagraph{Specific heat capacity in solids}

Experimentally, $C_V \sim aT + bT^3$ (electronic and lattice).
At high $T$, $C_V$ approaches constant.

\subparagraph{Dulong-Petit law}

Each atom has $3$ independent harmonic oscillators; each contributes $k_BT$ to
energy. Thus, $C_V = 3N_Ak_B = 3R$.
\begin{enumext}
  \item Only applicable to heavier elements (fail for light and
  strongly bonded atoms, such as diamond)
  \item Appliable only at high temperatures.
\end{enumext}

\paragraph{Einstein model}

Assume \emph{independent harmonic oscillators} with Einstein
frenquency $\omega_E$ and quantized
energy $E_n = \hbar\omega_E\ab(n + \frac12)$.

$N$ vibrating atoms need $3N$ independent 1D oscillators to
describe the motion.
The heat capacity can be obtained
\[
  C = \pdv{\bar E}T = -\pdv*[fun]{\frac1Z\pdv Z\beta}T
  = k_B(\beta\hbar\omega_E)^2
      \frac{\upe^{\beta\hbar\omega_E}}{(\upe^{\beta\hbar\omega_E} - 1)^2}
\]
\begin{enumext}
  \item High temperature ($\beta \to 0$), $C \to 3Nk_B$,
  approach to Dulong-Petit law at high temperature.
  \item Low temperature ($\beta \to \infty$), $C \to 0$,
  but exponentially not $T^3$.
\end{enumext}

\subparagraph{Limitation of Einstein model}

\begin{enumext}
  \item Assume all the vibrations share the same frequency.
  \item At low temperature, $\omega > k_BT/\hbar$ is frozen. The acoustic branch
  of the vibrations contribute a lot to the heat capacity.
\end{enumext}

\paragraph{Debye model}

\begin{enumext}
  \item Atoms are coupled to each other, there vibrations generate sound waves
  with a linear dispersion $\omega(k) = \nu k$.
  \item The sound waves are quantized similar to light waves, but there are
  three modes (one longitudinal two transverse) for each $\bm k$.
  \item There is a maximum oscillation frequency (Debye frequency), Debye
  temperature, and Debye radius
  \[
    \omega_D = (6\pi^2n)^{1/3} \nu,\quad
    T_D = \hbar\omega_D/k_B, \quad q_D = \omega_D/\nu.
  \]
  The heat capacity
  \[
    C = \pdv{\bar E}T
  = 9Nk_B\ab(\frac T{T_D})^3 \int_0^{T_D/T}
    \frac{x^4\upe^x}{(\upe^x - 1)^2} \d x.
  \]
  \begin{enumext}
    \item $T \gg T_D$: $C \to 3Nk_B$.
    \item $T \ll T_D$: $C \to \frac{12}{5}\pi^4 Nk_B T^3/T_D^3$.
  \end{enumext}
\end{enumext}

\subparagraph{Limitation of Debye model}

\begin{enumext}
  \item The cutoff seem like a cheat rather than physics.
  \item The assumption that sound waves exist even for very large values of $k$,
  but sound is based on a long-wavelength idea.
  \item $T_D$ is varying with temperature experimentally.
  \item Metal's heat capacity also have a term proportional
  to $T$.
\end{enumext}

\paragraph{Vibrations in a monatomic chain}

Atoms of mass $m$, spacing $a$, nearest-neighbor spring constant $\beta$.
EOM
\[
  m \ddot{u}_n = \beta(u_{n+1} + u_{n-1} - 2u_n).
\]
\begin{enumext}
  \item Wave solution gives dispersion $\omega = 2\sqrt{\beta/m} |\sin(ka/2)|$.
  \item For small $k$, $\omega \approx a\sqrt{\beta/m}|\bm k|$, looks
  like sound waves. Implies the velocity
  of sound (phase velocity) $v_s = v_p = \frac\omega k = a\sqrt{\frac\beta m}$.
  \item At the boundary of 1st BZ, $\omega_{\max} = 2\sqrt{\beta/m}$,
  $v_g = \pdv\omega k = 0$.
\end{enumext}

\subparagraph{Quantum modes: phonons}

Refer to a phonon as a Quantized lattice vibration.
\begin{enumext}
  \item For a normal mode of frequency $\omega$,
  $E_n = (n + \frac12)\hbar\omega$.
  \item Phonons are bosons, occupation given by Bose-Einstein distribution
  $n_B(\omega) = [e^{\hbar\omega/k_B T} - 1]^{-1}$.
\end{enumext}

\subparagraph{Crystal momentum}

Phonon carries crystal momentum $\hbar k$ (quasi-momentum, defined modulo
reciprocal lattice vectors). Conservation of crystal momentum governs
scattering processes.


\paragraph{Vibrations in a diatomic chain}

Two atoms per unit cell (masses $m$, $M$) with the dispersion relation
\[
  \omega_\pm^2
= \frac{\beta(M + m)}{Mm}\ab[1 \pm \sqrt{1 - \frac{4Mm}{(M + m)^2}\sin^2ak}],
\]
gives the acoustic ($\omega_-$) and optical ($\omega_+$) branches.
\begin{enumext}
  \item $k = 0$, $\omega_-^2 = 0$, $\omega_+^2 = 2\beta/\mu$.
  \item $k = \pi/(2a)$, $\omega_+^2 = 2\beta/m$, $\omega_-^2 = 2\beta/M$.
\end{enumext}
The gap at $\pm\frac\pi{2a}$ exists if $M \neq m$.
\begin{enumext}
  \item Acoustic mode: Atoms vibrate in same direction for $\omega_-$.
  \item Optical mode: Atoms vibrate in opposite direction for $\omega_-$,
  can couple to electromagnetic radiation (infrared).
\end{enumext}

\paragraph{Vibrations in 3D lattices}

For a unit cell with $n$ atoms, there are $3n$ branches
$3$ acoustic (one longitudinal, two transverse) and $3n - 3$ optical branches.

\paragraph{Experimental observations of vibrations}

X-ray scattering, neutron scattering, Raman spectroscopy, Brillouin
spectroscopy, etc., used to measure dispersion $\omega(k)$ and density of
states $g(\omega)$.

\begin{framed}\large\bfseries
  \#4 Crystal Symmetry and Structure
\end{framed}

\paragraph{Crystal definition}

Periodic, Basis, Symmetry, Anisotropy.

\subparagraph{Bravais lattice}

A lattice is repeating structure defined as an infinite set of points
\[
  \bm R = n_1\bm a_1 + n_2\bm a_2 + n_3\bm a_3,
\]
which can be constructed as sums of a set of linearly independent
primitive vectors $\bm a_i$. With the properties of
\emph{identical}, \emph{infinite extension},
and~\emph{contains crystal symmetry}.

\subparagraph{Primitive unit cell}

Contains only one lattice point, with
volume $\Omega = \bm a_1 \cdot (\bm a_2 \times \bm a_3)$.

\subparagraph{Wigner-Seitz cell}

Lines perpendicular and connecting the midpoints between a given lattice and all
nearby lattice sites.

\subparagraph{Crystallographic directions and planes}

Crystallographic directions $[l_1 l_2 l_3]$
represents $\bm n = l_1\bm a_1 + l_2\bm a_2 + l_3\bm a_3$.
A group of identical directions can be written as $\braket<l_1 l_2 l_3>$, like
\[
  \braket<100> = \{[100],\ [\bar100],\ [010],\ [0\bar10],\ [001],\ [00\bar1]\}.
\]
Planes: Miller indices $(hkl)$ obtained from reciprocals of intercepts
\[
  1/a:1/b:1/c = h:k:l
\]
with the corresponding family $\{hkl\}$.

\paragraph{Crystal symmetry}

\subparagraph{Symmetry operations}

\begin{enumext}
  \item $n$-fold rotation axis: An object that requires rotation of $2\pi/n$
  \item Reflection (mirror), Inversion, Rotoinversion.
\end{enumext}

\subparagraph{Crystal structures}

\begin{itemize}
  \item \ce{CsCl}: Simple cubic with basis \ce{Cs} at $[0,0,0]$,
  \ce{Cl} at $[\frac12,\frac12,\frac12]$.
  \item \ce{NaCl}: FCC with basis Na at $[0,0,0]$, Cl at $[\frac12,0,0]$.
  \item Diamond/Zinc blende (\ce{ZnS}): FCC with basis two \ce C
  (or \ce{Zn} and \ce{S}) atoms at $[0,0,0]$ and $[\frac14,\frac14,\frac14]$.
  \item Perovskite (\ce{ABO3}): Cubic with \ce A at corners, \ce B at body
  center, \ce O at face centers.
\end{itemize}

\subparagraph{Reciprocal lattice}

Defined by $\bm G \cdot \bm R = 2\pi m$. Primitive vectors
\[
  \bm b_i = \frac{2\pi (\bm a_j \times \bm a_k)}{\bm a_i \cdot (\bm a_j \times \bm a_k)}, \quad
  \bm a_i \cdot \bm b_j = 2\pi \delta_{ij}, \quad
  \Omega = \bm a_1 \cdot (\bm a_2 \times \bm a_3).
\]
Reciprocal lattice vectors are normal to real lattice planes $(hkl)$;
interplanar spacing $d_{hkl} = 2\pi/|\bm G_{hkl}|$. See \textbf{Problem 3.3}.

\subparagraph{Brillouin zone}

Wigner-Seitz unit cell of reciprocal lattice is the first BZ, contains all
wavevectors closer to origin than to any other reciprocal lattice point.
\begin{enumext}
  \item SC.
  $\Gamma:(0,0,0)$, $X:\pi/a(1,0,0)$, $R:\pi/a(1,1,1)$.
  \item FCC.
  $\Gamma:(0,0,0)$, $X:2\pi/a(1,0,0)$, \\
  $L:\pi/a(1,1,1)$, $R:2\pi/a(3/4,3/4,1)$.
  \item BCC.
  $\Gamma:(0,0,0)$, $X:2\pi/a(1,0,0)$, \\
  $P:\pi/a(1,1,1)$, $N:2\pi/a(1/2,1/2,0)$.
\end{enumext}
\begin{center}
  \includegraphics[width = .3\linewidth]{SC.png}
  \hfill
  \includegraphics[width = .3\linewidth]{FCC.png}
  \hfill
  \includegraphics[width = .3\linewidth]{BCC.png}
\end{center}

\subparagraph{Chemical bonding}

\[
  \text{Cohesive energy} = \text{energy of free atoms} - \text{crystal energy}.
\]

\begin{enumext}
  \item van der Waals (weak, induced dipole)
  \[
    U = 4\epsilon\ab[\ab(\frac\sigma R)^{12} - \ab(\frac\sigma R)^6].
  \]
  \item Ionic crystals (electrostatic, strong)
  \item Covalent (electron sharing, directional)
  \item Metallic (delocalized electrons)
  \item Hydrogen bonding (dipole-dipole).
\end{enumext}

\begin{framed}\large\bfseries
  \#5 Scattering Experiments
\end{framed}

\paragraph{X-Ray}

\begin{enumext}
  \item X-rays are electromagnetic waves. Atoms scatter X-ray waves
  elastically, primarily through the atoms' electrons.
  \item Some scattered waves can add constructively (constructive interference)
  in a few specific directions.
  \item The diffraction pattern of X-ray is highly relative to the crystal
  structure.
\end{enumext}

\subparagraph{Bragg's Law}

\[
  2d_{hkl} \sin\theta = n\lambda
\]

\subparagraph{Laue's condition}

Under elastic scattering
\[
  \bm k' - \bm k = \bm G.
\]

\subparagraph{Ewald construction for diffraction}

\begin{enumext}
  \item Draw a sphere in reciprocal space centered on the tip of the incident
  wave vector $\bm k$ of radius $k$ (so that it passed through the origin).
  \item There will be some wave vector satisfying the Laue condition if and
  only if some reciprocal lattice point (in addition to the origin lines on the
  surface of the sphere).
\end{enumext}

\subparagraph{Ewald construction}

\subparagraph{Geometrical structure factor}

\subparagraph{Diffraction methods}

See \textbf{Problem 4.2}.

\begin{framed}\large\bfseries
  \#6 Bloch Theorem
\end{framed}

\paragraph{Bloch theorem}
For electron in periodic potential $U(\bm r + \bm R) = U(\bm r)$, the wavefunction has form
\[
  \psi_{\bm k}(\bm r) = e^{i\bm k \cdot \bm r} u_{\bm k}(\bm r), \quad u_{\bm k}(\bm r + \bm R) = u_{\bm k}(\bm r).
\]
Crystal momentum $\hbar \bm k$; energy bands $E_n(\bm k)$ periodic in $\bm k$ with period $\bm G$.

\subparagraph{Brillouin zone}
Wavevector $\bm k$ can be restricted to first Brillouin zone. Number of allowed $\bm k$-states equals number of unit cells.

\subparagraph{Kronig–Penney model}
1D periodic square-well potential. Dispersion given by
\[
  \cos(kd) = \cos(\alpha a) \cosh(\beta b) - \frac{\alpha^2 - \beta^2}{2\alpha\beta} \sin(\alpha a) \sinh(\beta b),
\]
where $\alpha^2 = 2mE/\hbar^2$, $\beta^2 = 2m(V_0-E)/\hbar^2$, $d=a+b$. Solutions exist only for certain energy ranges (bands), separated by gaps.

\subparagraph{Failure of free electron theory}
Cannot explain insulators/semiconductors, temperature dependence of Lorenz number, sign anomalies in Hall effect, etc.

\subparagraph{de Broglie hypothesis}
Particles have wave nature; $\lambda = h/p$. Confirmed by electron diffraction.

\begin{framed}\large\bfseries
  \#7 Band Theory I-NFE and TB model
\end{framed}

\paragraph{Nearly free electron (NFE) model}
Weak periodic potential treated as perturbation. Unperturbed states are plane waves. Degenerate perturbation at Brillouin zone boundaries opens energy gaps. For 1D, gap magnitude $E_g = 2|U_n|$, where $U_n$ is Fourier component of potential.

\subparagraph{Tight binding (TB) model}
Electrons tightly bound to atoms; wavefunction = linear combination of atomic orbitals (LCAO):
\[
  \psi(\bm r) = \sum_m c_m \varphi_i(\bm r - \bm R_m).
\]
Energy bands:
\[
  E(\bm k) = \varepsilon_i - \sum_s J(\bm R_s) e^{-i\bm k \cdot \bm R_s},
\]
where $J(\bm R_s)$ is hopping integral. Nearest-neighbor approximation often used.

\subparagraph{Examples}
Simple cubic with s-orbitals:
\[
  E(\bm k) = \varepsilon_s - J_0 - 2J_1 (\cos k_x a + \cos k_y a + \cos k_z a).
\]
Bandwidth proportional to $J_1$. Effective mass near band bottom: $m^* = \hbar^2/(2J_1 a^2)$.

\subparagraph{Comparison}
NFE: good for nearly free electrons (conduction bands). TB: good for tightly bound electrons (valence bands).

\begin{framed}\large\bfseries
  \#8 Band Theory II
\end{framed}

\paragraph{Intrinsic semiconductors}
Pure semiconductor; carrier concentrations $n = p = n_i$. At finite $T$,
\[
  n = N_c e^{-(E_c - E_f)/k_B T}, \quad p = N_v e^{-(E_f - E_v)/k_B T},
\]
where $N_c, N_v$ are effective densities of states. Intrinsic carrier concentration
\[
  n_i = \sqrt{N_c N_v} e^{-E_g/(2k_B T)}.
\]
Fermi level near midgap: $E_f = E_v + \frac{1}{2}E_g - \frac{3}{4}k_B T \ln(m_e^*/m_h^*)$.

\subparagraph{Doped semiconductors}
\begin{itemize}
  \item n-type: donors (group V) add electrons; $n \approx N_d$, $p = n_i^2/N_d$.
  \item p-type: acceptors (group III) add holes; $p \approx N_a$, $n = n_i^2/N_a$.
\end{itemize}
Fermi level shifts toward donor/acceptor levels.

\subparagraph{Carrier scattering}
Impurity scattering: mobility $\propto T^{3/2}/N_i$. Lattice scattering: mobility $\propto T^{-3/2}$ (acoustic phonons) or $T^{-1/2}$ (optical phonons). Total mobility: $1/\mu_{\text{total}} = \sum 1/\mu_i$.

\subparagraph{p-n junction}
Formed by joining p- and n-type semiconductors. Built-in potential $V_0 = \frac{k_B T}{e} \ln(N_a N_d / n_i^2)$. Depletion width $W = \sqrt{\frac{2\varepsilon (N_a+N_d) V_0}{e N_a N_d}}$.
Current–voltage characteristic (ideal diode):
\[
  I = I_0 (e^{eV/(k_B T)} - 1).
\]
Reverse breakdown: avalanche (impact ionization) or Zener (tunneling).

\subparagraph{Devices}
\begin{itemize}
  \item LED: forward-biased p-n junction; recombination emits light.
  \item Solar cell: photon absorption creates electron-hole pairs, separated by built-in field.
  \item BJT: npn or pnp structure; amplifies current.
  \item MOSFET: metal-oxide-semiconductor field-effect transistor; voltage-controlled channel.
  \item Schottky barrier: metal-semiconductor contact; rectifying or ohmic.
\end{itemize}

\subparagraph{Heterojunctions}
Different semiconductors joined; band alignments: type I (straddling), type II (staggered), type III (broken gap). Used in 2D electron gases, optical devices, tunneling diodes.

\begin{framed}\large\bfseries
  \#9 Physics of Semiconductors
\end{framed}

\paragraph{Intrinsic semiconductors}
Pure semiconductor; carrier concentrations $n = p = n_i$. At finite $T$,
\[
  n = N_c e^{-(E_c - E_f)/k_B T}, \quad p = N_v e^{-(E_f - E_v)/k_B T},
\]
where $N_c, N_v$ are effective densities of states. Intrinsic carrier concentration
\[
  n_i = \sqrt{N_c N_v} e^{-E_g/(2k_B T)}.
\]
Fermi level near midgap: $E_f = E_v + \frac{1}{2}E_g - \frac{3}{4}k_B T \ln(m_e^*/m_h^*)$.

\subparagraph{Doped semiconductors}
\begin{itemize}
  \item n-type: donors (group V) add electrons; $n \approx N_d$, $p = n_i^2/N_d$.
  \item p-type: acceptors (group III) add holes; $p \approx N_a$, $n = n_i^2/N_a$.
\end{itemize}
Fermi level shifts toward donor/acceptor levels.

\subparagraph{Carrier scattering}
Impurity scattering: mobility $\propto T^{3/2}/N_i$. Lattice scattering: mobility $\propto T^{-3/2}$ (acoustic phonons) or $T^{-1/2}$ (optical phonons). Total mobility: $1/\mu_{\text{total}} = \sum 1/\mu_i$.

\subparagraph{p-n junction}
Formed by joining p- and n-type semiconductors. Built-in potential $V_0 = \frac{k_B T}{e} \ln(N_a N_d / n_i^2)$. Depletion width $W = \sqrt{\frac{2\varepsilon (N_a+N_d) V_0}{e N_a N_d}}$.
Current–voltage characteristic (ideal diode):
\[
  I = I_0 (e^{eV/(k_B T)} - 1).
\]
Reverse breakdown: avalanche (impact ionization) or Zener (tunneling).

\subparagraph{Devices}
\begin{itemize}
  \item LED: forward-biased p-n junction; recombination emits light.
  \item Solar cell: photon absorption creates electron-hole pairs, separated by built-in field.
  \item BJT: npn or pnp structure; amplifies current.
  \item MOSFET: metal-oxide-semiconductor field-effect transistor; voltage-controlled channel.
  \item Schottky barrier: metal-semiconductor contact; rectifying or ohmic.
\end{itemize}

\subparagraph{Heterojunctions}
Different semiconductors joined; band alignments: type I (straddling), type II (staggered), type III (broken gap). Used in 2D electron gases, optical devices, tunneling diodes.

\begin{framed}\large\bfseries
  \#10 Electron in a Magnetic Field
\end{framed}

\paragraph{Bohr–Sommerfeld quantization}\leavevmode

\subparagraph{Quantization rule}\leavevmode

\begin{enumext}
  \item The number of wavelengths along a closed trajectory must be integer:
    \[
      \frac{1}{\hbar} \oint \bm p \cdot \d\bm r = 2\pi N, \quad N \in \mathbb Z.
    \]
  \item In a magnetic field \(\bm B = \nabla \times \bm A\), the canonical momentum becomes \(\bm p \rightarrow \bm p - e\bm A\).
  \item The quantization condition becomes
    \[
      \frac{1}{\hbar} \oint (\bm p - e\bm A) \cdot \d\bm r = 2\pi N.
    \]
  \item This leads to quantization of the magnetic flux enclosed by the orbit:
    \[
      \phi = N \phi_0, \quad \phi_0 = \frac{h}{e}.
    \]
\end{enumext}

\subparagraph{Onsager quantization condition}\leavevmode

\begin{enumext}
  \item In \(k\)-space, the area enclosed by an orbit is quantized:
    \[
      S_k = \frac{2\pi eB}{\hbar} \ab(N + \frac12 - \frac{\gamma}{2\pi}),
    \]
    where \(\gamma\) is the Berry phase.
  \item For free electrons (\(S_k = \pi k^2\)), the energy becomes quantized:
    \[
      E_N = \hbar \omega_c \ab(N + \frac12 - \frac{\gamma}{2\pi}), \quad \omega_c = \frac{eB}{m^*}.
    \]
  \item Without Berry phase (\(\gamma = 0\)):
    \[
      E_N = \hbar \omega_c \ab(N + \frac12).
    \]
\end{enumext}

\paragraph{Landau quantization}\leavevmode

\subparagraph{Quantum mechanical solution}\leavevmode

\begin{enumext}
  \item Hamiltonian for an electron in a magnetic field (Landau gauge \(\bm A = (0, Bx, 0)\)):
    \[
      \hat H = \frac{1}{2m^*} (\bm p - e\bm A)^2.
    \]
  \item The motion parallel to \(\bm B\) is free, while the perpendicular motion is quantized into Landau levels (LLs):
    \[
      E = \frac{\hbar^2 k_z^2}{2m^*} + \ab(N + \frac12) \hbar \omega_c, \quad N = 0,1,2,\dots
    \]
  \item In 2D (\(k_z=0\)): \(E_N = \ab(N + \frac12) \hbar \omega_c\).
\end{enumext}

\subparagraph{Degeneracy of Landau levels}\leavevmode

\begin{enumext}
  \item Each Landau level has a degeneracy per unit area:
    \[
      n_B = \frac{eB}{h} \quad (\text{per spin}).
    \]
  \item With spin degeneracy \(g_s = 2\): \(n_B = 2\frac{eB}{h}\).
  \item The filling factor \(\nu = \frac{n_e}{n_B} = \frac{hn_e}{eB}\).
\end{enumext}

\subparagraph{Density of states}\leavevmode

\begin{enumext}
  \item In 3D, the DOS becomes a series of peaks (Landau tubes) broadened by scattering.
  \item In 2D, the DOS is a set of delta functions at each Landau level, broadened by disorder.
  \item The broadening \(\Delta E \sim \hbar / \tau_q\), where \(\tau_q\) is the quantum lifetime.
  \item Landau levels are resolved when \(\omega_c \tau_q \gg 1\).
\end{enumext}

\subparagraph{Zeeman splitting}\leavevmode

\begin{enumext}
  \item Zeeman term adds spin splitting:
    \[
      E = \ab(N + \frac12) \hbar \omega_c \pm \frac12 g^* \mu_B B.
    \]
  \item The spin susceptibility \(\chi \propto g^* m^*\).
\end{enumext}

\paragraph{Shubnikov–de Haas (SdH) oscillations}\leavevmode

\subparagraph{Basic phenomenology}\leavevmode

\begin{enumext}
  \item Oscillations in longitudinal resistivity \(\rho_{xx}\) as a function of \(1/B\).
  \item Periodicity in \(1/B\):
    \[
      \Delta\ab(\frac1B) = \frac{2e}{h n_e}.
    \]
  \item The oscillation frequency \(B_F = \frac{h}{2\pi e} S_k\) gives the extremal cross-sectional area of the Fermi surface.
\end{enumext}

\subparagraph{Lifshitz–Kosevich (LK) formalism}\leavevmode

\begin{enumext}
  \item Amplitude of SdH oscillations:
    \[
      \frac{\Delta\rho_{xx}}{\rho_0} = R_T R_D R_S \cos\ab[2\pi\ab(\frac{B_F}{B} + \frac12 + \beta)],
    \]
    where
    \begin{align*}
      R_T &= \frac{2\pi^2 k_B T / \hbar\omega_c}{\sinh(2\pi^2 k_B T / \hbar\omega_c)} \quad \text{(thermal damping)},\\
      R_D &= \upe^{-\pi/\omega_c \tau_q} \quad \text{(Dingle factor)},\\
      R_S &= \cos\ab(\frac{\pi g^* m^*}{2 m_0}) \quad \text{(spin damping)}.
    \end{align*}
  \item Phase factor \(\beta = \gamma/2\pi\) contains the Berry phase \(\gamma\).
\end{enumext}

\subparagraph{Extracting parameters}\leavevmode

\begin{enumext}
  \item From temperature dependence of the amplitude → cyclotron mass \(m^*\).
  \item From Dingle plot → quantum lifetime \(\tau_q\).
  \item From phase \(\beta\) → Berry phase (\(\beta = 1/2\) for Dirac fermions).
\end{enumext}

\paragraph{Integer quantum Hall effect (IQHE)}\leavevmode

\subparagraph{Basic observations}\leavevmode

\begin{enumext}
  \item At high magnetic fields (\(\omega_c \tau \gg 1\)), Hall resistivity \(\rho_{xy}\) forms plateaus at quantized values:
    \[
      \rho_{xy} = \frac{h}{\nu e^2}, \quad \nu = 1,2,3,\dots
    \]
  \item Simultaneously, longitudinal resistivity \(\rho_{xx} \to 0\).
  \item Quantized Hall conductance \(\sigma_{xy} = \nu \frac{e^2}{h}\).
\end{enumext}

\subparagraph{Edge states}\leavevmode

\begin{enumext}
  \item In a finite sample, confinement potential creates edge channels.
  \item Edge states are chiral (unidirectional) and carry the current.
  \item Backscattering is suppressed → dissipationless transport.
\end{enumext}

\subparagraph{Localization and plateaus}\leavevmode

\begin{enumext}
  \item Disorder localizes states in the tails of Landau levels.
  \item Extended states exist only near the center of each Landau level.
  \item As the Fermi level moves through localized states, \(\sigma_{xx}\) remains zero → Hall plateaus.
\end{enumext}

\paragraph{Fractional quantum Hall effect (FQHE)}\leavevmode

\subparagraph{Basic observations}\leavevmode

\begin{enumext}
  \item Occurs in very clean 2DEG at high \(B\) and low \(T\).
  \item Hall plateaus at fractional filling factors \(\nu = p/q\) (e.g., \(1/3, 2/5, 3/7\)).
  \item \(\rho_{xx}\) shows minima at the same fillings.
\end{enumext}

\subparagraph{Composite fermion picture}\leavevmode

\begin{enumext}
  \item Each electron binds an even number \(2j\) of magnetic flux quanta → composite fermion (CF).
  \item CFs experience an effective magnetic field \(B^* = B - 2j \phi_0 n_e\).
  \item The FQHE of electrons is the IQHE of CFs.
  \item Principal Jain sequence: \(\nu = \frac{j}{2j \pm 1}\).
\end{enumext}

\subparagraph{Wigner crystal}\leavevmode

\begin{enumext}
  \item At very low density and high \(B\), electrons can crystallize into a Wigner crystal.
  \item Characterized by vanishing conductivity and pinning mode resonances.
\end{enumext}

\begin{framed}\large\bfseries
  \#11, \#12 Dirac Band Theory
\end{framed}

\paragraph{Band structure of monolayer graphene}\leavevmode

See \textbf{Problem 9.1}. $E = \epsilon_0 \pm |f(\bm k)|$,
$f(\bm k) = \sum_{i=1}^3 \upe^{\iu\bm k \cdot \delta_l}$.

\subparagraph{More general expression}

Consider specific example of two orbitals per unit cell
\[
  E_\pm = \frac{\epsilon_0 \pm \gamma_0|f(\bm k)|}{1 \mp s_0|f(\bm k)|}
\]
$E_+$ and $E_-$ are the conduction band and valence band.
\begin{enumext}
  \item Two bands: conduction ($\pi$) and valence ($\pi^*$) band.
  \item No energy gap at the $K$, $K' = \frac{4\pi}{3\sqrt3a}(\pm1,0)$ points.
  \item The reasonable case: $\epsilon_0 = 0$, $s_0 = 0$,
  $E_\pm = \pm\gamma_0|f(\bm k)|$. Then at $K$, $K'$ points $E_\pm = 0$.
  Indicating that there is no coupling between the A and B sublattices.
  They support the same quantum states, leading to a degeneracy point.
  \item $K$, $K'$ points:
  $4\text{-fold degeneracy} = 2\ \text{sublattices} + 2\ \text{valleys}$.
\end{enumext}

\subparagraph{Linear expansion near $K$, $K'$ points}\leavevmode

See \textbf{Problem 9.2}: Resemble massless particle (photons) $E = cp$,
but with 300 times slower speed.

\subparagraph{Pseudospin degree of freedom}

Pseudospin is a degree of freedom related to the relative amplitude of Bloch
function on the A or B sublattice.

\subparagraph{Chirality}

The orientation of the pseudospin is related to the direction of the momentum,
The amplitudes on the A or B sublattice of eigenstate depend
on $\theta_q = \arctan p_y/p_x$.

Chirality is reversed in the $K'$ valley.
\begin{enumext}
  \item Conduction band electron: $\bm \sigma \cdot \bm n = 1$.
  \item Valence band electron: $\bm \sigma \cdot \bm n = -1$.
\end{enumext}
Chirality conserved: no backscattering between Dirac cones.

\subparagraph{Klein tunneling}

If a given potential doesn't break A-B symmetry, then it is unable to
influence the pseudospin degree of freedom, therefore, be conserved upon
scattering.

Klein paradox:
\begin{enumext}
  \item Transmission of relativistic particles unimpeded even by
  highest barriers.
  \item Physical picture: particle/hole conversion.
\end{enumext}

\subparagraph{Berry phase}

The Berry phase is the same for conduction and valley band.
Berry phase around $K$, $K'$ point $\phi_B = \pm\pi$.

\subparagraph{Trigonal warping beyond linearization} \leavevmode

Consider the next nearest-neighbor hopping $t'$.
\[
  E(\bm q) \approx 3t' \pm v_F|\bm q|
- \ab(\frac{9t'a^2}{4} \pm \frac{3t}{8}a^2\sin3\arctan\frac{q_x}{q_y})\bm q^2
\]
Trigonal warping
\begin{enumext}
  \item $t'$ shifts in energy and breaks electron-hole symmetry.
  \item The second order part depends on the direction in momentum space and
  has a threefold symmetry.
\end{enumext}

\subparagraph{Boron nitride: Dirac equation with a mass}



\paragraph{Band structure of bilayer graphene}\leavevmode

\paragraph{Landau levels in graphene}\leavevmode


\hrule

\paragraph{Introduction to graphene}\leavevmode

\subparagraph{Basic structure}\leavevmode

\begin{enumext}
  \item 2D honeycomb lattice of carbon atoms, sp\(^2\) hybridization.
  \item Two sublattices (A and B) per unit cell.
  \item Lattice constant \(a \approx 2.46\,\text{\AA}\).
\end{enumext}

\subparagraph{Electronic properties}\leavevmode

\begin{enumext}
  \item Zero‑gap semiconductor (semimetal) with linear dispersion at \(K\) and \(K'\) points.
  \item Charge carriers behave as massless Dirac fermions with Fermi velocity \(v_F \approx 10^6\,\text{m/s}\).
  \item High mobility, excellent thermal conductivity, optically transparent (\(\pi\alpha \approx 2.3\%\) absorption).
\end{enumext}

\paragraph{Band structure of monolayer graphene}\leavevmode

\subparagraph{Tight‑binding model}\leavevmode

\begin{enumext}
  \item Nearest‑neighbor hopping integral \(\gamma_0 \approx 3.0\,\text{eV}\).
  \item Hamiltonian in basis \(( \psi_A, \psi_B )\):
    \[
      H = \begin{pmatrix}
        0 & -\gamma_0 f(\bm k)\\
        -\gamma_0 f^*(\bm k) & 0
      \end{pmatrix},
      \quad
      f(\bm k) = \sum_{i=1}^3 \upe^{\iu \bm k \cdot \bm\delta_i}.
    \]
  \item Energy dispersion:
    \[
      E(\bm k) = \pm \gamma_0 |f(\bm k)|.
    \]
  \item At the \(K\) points, \(f(\bm k)=0\) → Dirac cones.
\end{enumext}

\subparagraph{Low‑energy effective Hamiltonian}\leavevmode

\begin{enumext}
  \item Near \(K\) (\(\xi=+1\)) and \(K'\) (\(\xi=-1\)):
    \[
      H_\xi(\bm p) = v_F \begin{pmatrix}
        0 & \xi p_x - \iu p_y\\
        \xi p_x + \iu p_y & 0
      \end{pmatrix}
      = v_F (\xi \sigma_x p_x + \sigma_y p_y).
    \]
  \item Dispersion: \(E_{\pm}(\bm p) = \pm v_F |\bm p|\).
  \item Eigenstates (pseudospinors):
    \[
      \psi_{\pm,\xi} = \frac{1}{\sqrt2} \begin{pmatrix}
        1\\
        \pm \xi \upe^{\iu \xi \theta_{\bm p}}
      \end{pmatrix}, \quad
      \theta_{\bm p} = \arctan(p_y/p_x).
    \]
\end{enumext}

\subparagraph{Berry phase and chirality}\leavevmode

\begin{enumext}
  \item Berry phase \(\gamma = \pi\) for a closed loop around a Dirac point.
  \item Chirality (helicity) operator: \(\hat h = \bm\sigma \cdot \bm n\), where \(\bm n = \bm p/|\bm p|\).
  \item Chirality is conserved → no back‑scattering → Klein tunneling.
\end{enumext}

\subparagraph{Density of states}\leavevmode

\begin{enumext}
  \item Linear DOS near Dirac point:
    \[
      g(E) = \frac{2|E|}{\pi \hbar^2 v_F^2}.
    \]
  \item Carrier density \(n\) and Fermi energy:
    \[
      E_F = \hbar v_F \sqrt{\pi |n|}, \quad
      k_F = \sqrt{\pi |n|}.
    \]
\end{enumext}

\paragraph{Band structure of bilayer graphene}\leavevmode

\subparagraph{Tight‑binding model}\leavevmode

\begin{enumext}
  \item Bernal (AB) stacking: interlayer coupling \(\gamma_1 \approx 0.39\,\text{eV}\) between A2 and B1.
  \item Hamiltonian in basis (A1, B1, A2, B2):
    \[
      H = \begin{pmatrix}
        0 & -\gamma_0 f & 0 & 0\\
        -\gamma_0 f^* & 0 & \gamma_1 & 0\\
        0 & \gamma_1 & 0 & -\gamma_0 f\\
        0 & 0 & -\gamma_0 f^* & 0
      \end{pmatrix}.
    \]
  \item Four bands: two high‑energy bands split by \(\pm\gamma_1\), and two low‑energy bands touching at \(K\).
\end{enumext}

\subparagraph{Low‑energy effective Hamiltonian}\leavevmode

\begin{enumext}
  \item For the low‑energy bands, effective two‑component Hamiltonian:
    \[
      H_{2,\xi}(\bm p) = -\frac{1}{2m^*} \begin{pmatrix}
        0 & (\xi p_x - \iu p_y)^2\\
        (\xi p_x + \iu p_y)^2 & 0
      \end{pmatrix},
      \quad m^* = \frac{\gamma_1}{2v_F^2}.
    \]
  \item Parabolic dispersion: \(E_{\pm}(\bm p) = \pm \frac{p^2}{2m^*}\).
  \item Berry phase \(\gamma = 2\pi\).
\end{enumext}

\subparagraph{Trigonal warping}\leavevmode

\begin{enumext}
  \item Due to skew interlayer hopping \(\gamma_3\), the low‑energy dispersion becomes anisotropic.
  \item At very low energies, the Fermi surface splits into four Dirac pockets.
\end{enumext}

\subparagraph{Electrically tunable band gap}\leavevmode

\begin{enumext}
  \item A perpendicular electric field breaks inversion symmetry → opens a band gap.
  \item Gap size \(\Delta \propto\) applied field, up to \(\sim \gamma_1\).
\end{enumext}

\paragraph{Landau levels in graphene}\leavevmode

\subparagraph{Monolayer graphene}\leavevmode

\begin{enumext}
  \item Landau level energies:
    \[
      E_N = \pm \hbar \omega_c \sqrt{2N}, \quad
      \omega_c = \frac{v_F}{l_B}, \quad
      l_B = \sqrt{\frac{\hbar}{eB}}.
    \]
  \item Zero‑energy Landau level (\(N=0\)) is half‑filled in neutral graphene.
  \item Fourfold degeneracy (spin × valley).
\end{enumext}

\subparagraph{Bilayer graphene}\leavevmode

\begin{enumext}
  \item Landau level energies:
    \[
      E_N = \pm \hbar \omega_c \sqrt{N(N-1)}, \quad
      \omega_c = \frac{eB}{m^*}.
    \]
  \item Two zero‑energy levels (\(N=0,1\)) → eightfold degeneracy.
\end{enumext}

\subparagraph{Quantum Hall effect}\leavevmode

\begin{enumext}
  \item Monolayer: Hall conductivity \(\sigma_{xy} = \pm (4N+2) \frac{e^2}{h}\) (half‑integer QHE).
  \item Bilayer: \(\sigma_{xy} = \pm 4N \frac{e^2}{h}\) (integer QHE).
  \item Room‑temperature QHE observable due to large Landau level spacing.
\end{enumext}